\section{Introduction}
\label{s:intro}
             

The topic of the paper is the comparison between 
\emph{Operational Game semantics} (OGS) and the \emph{$\pi$-calculus}, as generic
models or frameworks for the semantics of higher-order languages.  

Game semantics~\cite{10.1006/inco.2000.2917,10.1006/inco.2000.2930} provides
intensional models of higher-order languages,
where the denotation of a program brings up its possible  interactions with the
surrounding context.
Distinct points of game semantics are the 
rich categorical structure and the emphasis on compositionality.
Game semantics provides a modular characterization of higher-order languages with
 computational effects like control operators~\cite{10.5555/788019.788859}, 
 mutable store~\cite{abramsky1997linearity,10.5555/788020.788891} 
 or concurrency~\cite{ghica2004angelic,10.1007/11944836_38}.
 This gives rise to the ``Semantic Cube''~\cite{10.5555/645726.667210},
 a characterization of the absence of such computational effects
 in terms of appropriate restrictions on the interactions, with conditions
 like \emph{alternation, well-bracketing, visibility} or \emph{innocence}.
For instance, well-bracketing corresponds to the absence of control operators
like $\callcc$.

Game semantics has spurred Operational Game Semantics
(OGS)~\cite{jagadeesan2009open,jeffrey2005java,laird2007,10.1016/j.entcs.2012.08.013,10.1145/2603088.2603150},
as a way to describe the interactions of a program with its 
environment by embedding  programs into appropriate configurations and then defining rules
that turn such configurations into an LTS.
Besides minor differences in the representation of causality between actions,
the main distinction with ``standard'' game semantics
is in the way in which the denotation of programs is obtained:
via an LTS, rather than, 
compositionally, by induction on the structure of the programs (or their types).
It is nonetheless possible to establish a formal correspondence between these two 
representations~\cite{10.1145/2603088.2603150}.

OGS is particularly effective on higher-order programs.
To avoid being too intensional, functional values exchanged between the program and its environment 
are represented as atoms, seen as free variables.  
Therefore OGS configurations include open terms. 
The basic actions in the LTS produced by OGS represent
the calls and returns of functions between
a program
and its environment.
% The denotation of a program is then the  behaviour (usually seen as a set of traces)  generated by its associated LTS.
% The basic actions in such behaviours are the call and returns of functions between
%  the program
% and its environment.
% % OGS can be seen as an operational semantics or even as  an abstract machine for 
% % representing interactions of (possibly open) terms. 
% % as a sequence of actions, i.e. a trace.
The OGS semantics
has been shown  fully-abstract, that is, to characterize observational equivalence,
for a wide class of
programming languages, including effectful subsets of 
ML~\cite{laird2007,10.1145/2933575.2934509},
fragments of Java~\cite{jeffrey2005java}, aspect-oriented 
programs~\cite{jagadeesan2009open}. 
The conditions  in the above-mentioned Semantic Cube (alternation,
well-bracketing, etc.)
 equally apply to OGS.

In this paper, we consider forms of OGS
for the  pure untyped call-by-value  $\lambda$-calculus,
which enforce some of such conditions.
Specifically, we consider:
\begin{itemize}
 \item an \emph{Alternating} OGS, where only one term can be run at a time, and 
the control of the interactions alternates between the  term and the environment;
 \item a \emph{Concurrent} OGS, where multiple terms can be run in parallel;
\item a \emph{Well-Bracketed} OGS, where the calling convention between a term and its
environment follows a stack discipline.
\dsOLD{mention also complete traces?}
\end{itemize}


\vskip .2cm 
\noindent
The $\pi$-calculus is the paradigmatical name-passing calculus, that is, a calculus where
names (a synonym for `channels') may be passed around.  In the literature about the
$\pi$-calculus, and more generally in  Programming Language theory,
Milner's work on functions as processes~\cite{Mil92s}, which
shows how the evaluation strategies of {\em call-by-name $\lambda$-calculus} and {\em
  call-by-value $\lambda$-calculus}~\cite{Abr88,Plo75} can be
faithfully mimicked, is generally considered a landmark.  
% The
% comparison with the $\lambda$-calculus is a significant expressiveness test for the
% $\pi$-calculus.  
The work promotes the $\pi$-calculus to be a model for
 higher-order programs, and  
provides the means to study $\lambda$-terms in contexts other than
the purely sequential ones and with the instruments available to reason about 
processes.  
In the paper, $\pi$-calculus is actually meant to be 
the  \emph{Internal $\pi$-calculus} (\piI), 
a subset of the original $\pi$-calculus 
in which only fresh names may be exchanged among processes \cite{San95i}.
The use of \piI 
avoids  a few shortcomings
of Milner's encodings, notably for call-by-value; e.g., 
the  failure of the $\beta_v$ rule (i.e., the encodings of $ (\lambda x.M ) V$ and $M
\sub V x$ may be behaviourally distinguishable in  $ \pi$).
% ,  and    
% other desirable call-by-value equalities such as  $  I(xV)  = x V
% $, where $I$ is the identity $\lambda$-term.
%  
% Milner's original encoding of call-by-value $\lambda$-calculus had a few shortcomings \cite{Durier}: 
% it fails to validate the $\beta_v$ rule (i.e., the encodings of $ (\lambda x. ) V$ and $M
% \sub V x$ may be behaviourally distinguishable in  $ \pi$); and    
% desirable call-by-value equalities  fail, an example being $  I(xV)  = x V
% $, where $I$ is the identity $\lambda$-term).  
% A way of  overcoming these problems is to tranport the encoding into a subcalculus of the
% $\pi$-calculus, called \emph{Internal $\pi$-calculus} and indicated as \piI, 
% in which  any name emitted in an  output is fresh \cite{San95i}. Another advantage of \piI
% is  a simpler
%  algebraic and operational theory.  

Further investigations into Milner's  encodings \cite{San93d,San95lazy,DurierHS18} have revealed 
what is the equivalence induced on 
 $\lambda$-terms by the encodings, whereby
two $\lambda$-terms are equal if their  encodings are behaviourally
equivalent (i.e., bisimilar) $\piI$ terms. 
In call-by-value, this equivalence is \emph{eager normal-form bisimilarity}
\cite{lassentrees}, 
a tree structure
proposed  by Lassen 
(and indeed sometimes referred to as `Lassen's trees')
as the call-by-value  counterpart 
   of
B{\"o}hm Trees (or L{\'e}vy-Longo Trees). 

In a nutshell, when  used to give semantics to a language, 
major strengths of the $\pi$-calculus are its algebraic structure and the related algebraic
properties and proof techniques;
major strengths of     OGS   are its proximity to the source
language~--- the configurations of OGS are built directly from the terms of the source
language,  as opposed to an encoding as in the $\pi$-calculus~--- 
and its flexibility~--- the semantics can be tuned to account for specific features of the
source language like control operators or references. 

The general goal of this paper is to show that there is a tight and precise correspondence
between OGS and $\pi$-calculus as models of programming languages, and that such a
correspondence may be profitably used to take advantage of the strengths of the two models. 
We carry out the above program in the specific case of (untyped) call-by-value
$\lambda$-calculus,  $\LasV$, 
which is richer and (as partly suggested above) with some more subtle
aspects than call-by-name.
We also sketch at the end of the paper a similar correspondence for the
untyped call-by-name $\lambda$-calculus.
Analogies and similarities between game semantics and  $\pi$-calculus have  been
pointed  out in various papers in the literature (e.g., \cite{10.1145/224164.224189,berger2001sequentiality}; see
Section~\ref{s:cf}), and used to, e.g., explain game semantics using $\pi$-like processes,
and enhance type systems for $\pi$-terms.  In this paper, in contrast, we carry out a direct comparison
between the two models, on their interpretation of functions.    


We take the (arguably) canonical representations of $\LasV$  into \piI 
% (presented in Section~\ref{s:encoLpi})
and OGS.
% (in Section~\ref{s:cogs}).
The latter representation is Milner's encoding,
rewritten in \piI.  We consider two variant  behaviours for the \piI terms, respectively
produced by  the ordinary LTS
of \piI, and by  an `output-prioritised' LTS, \opLTS, 
% (Section~\ref{subsec:seq-ts-pi}) 
in which   input actions
may be observed only in the absence of outputs and internal actions. 
Intuitively, the \opLTS is intended to respect sequentiality  constraints in the
%observed 
\piI terms:  
an  output action stands for an ongoing computation (for instance, returning  the result of
 a previous request) whereas 
an  input
 action starts a 
   new computation
 (for instance, a request of a certain service); therefore,  
in a sequential system, an output action should have priority over input actions.  
 % output actions represent ongoing computations (for instance, returning  the result of
 % a previous interrogation of the term) and, in a sequential sytstem,    must therefore be consumed before any input
 % actions,   as  input actions represent 
 %   new computations
 % of the system (for instance, requesting a service with certain parameters). 
 For OGS, the $\LasV$ representation is the straightforward adaptation of the OGS
representations of typed $\lambda$-calculi in the literature, e.g., \cite{laird2007}
designed to build fully-abstract models for languages with mutable store.
As outlined above, a few variants exist:
% there are various possible constraints 
% \ds{``constraints''? }
%on the definition of the OGS;
% a few variants of the representation are possible;
%, depending on the features of OGS that one  wishes to emphasize; 
we consider the  alternating, concurrent,
and well-bracketed variants.

We also consider the refinements obtained by focusing on  \emph{complete traces}, 
intuitively, traces that describe ``terminating'' interactions.
Complete traces have been introduced in~\cite{abramsky1997linearity} to obtain an
effective presentation of a fully abstract model for a language with higher-order store.

We then develop a  thorough comparison between the behaviours of the OGS and \piI
representations. 
For this, we define a mapping from OGS configurations to \piI processes. We also exploit
the fact that, syntactically, the actions in the  OGS and   \piI LTSs are the same.
We derive a tight correspondence between the two models, which allows us to transfer techniques
and  to 
switch  freely between the two models in the analysis of 
the OGS and \piI representations of $\LasV$, 
so to establish new results or obtain new proofs.
On these aspects, our main results  are the following: 

\begin{enumerate}
\item
We show that
the representation of $\LasV$ in the Alternating 
  OGS  is behaviourally the same as the representation in \piI assuming the \opLTS.
Thus the semantics on $\lambda$-terms induced by the OGS and \piI representations
coincide. 
%  By `behaviourally the same' we mean that each
% $\lambda$-term  yields   behaviours  in OGS and \piI  that are bisimilar. 
% As a consequence, the equivalence induced on $\LasV$ by the OGS and $\piI$
% representations is the same. 
 The same results are obtained between the Concurrent  OGS and \piI under  its ordinary
 LTS. 
% \item  The equivalences produced in (1) and (2)  coincide, and coincide with Lassen's
%   trees. 
% We  prove  the same results for call-by-name $\lambda$-calculus, $\LasN$, and 
% the     L{\'e}vy-Longo Tree equality.


\item 
We   transfer 
 `bisimulation up-to techniques' for  $\piI$, notably a form of `up-to context',
 onto  (Concurrent) OGS.
The result is a powerful technique, called `up-to composition', that allows us to split an OGS
configuration into more elementary configurations during the bisimulation game. 
% %The technique was inspired by the mapping of OGS configurations into \piI and
% Its soundness is derived via the mapping   of OGS configurations into \piI. 


\item
We show that the semantics induced on $\LasV$ by the Alternating 
and by the Concurrent OGS
are the same, both when the equality in OGS is  based  on traces and
when it is based on bisimulation. 
In other words, all the OGS views of $\LasV$
(Alternating or Concurrent, traced-based or bisimulation-based) coincide. 
Moreover,  we show 
that 
 such
induced semantics is the equality of 
Lassen's trees. We derive the result in two ways:
one  in which we  directly import it from \piI; the other
%, presented in Appendix~\ref{s:enf},
in which we lift eager normal-form bisimulations into  OGS bisimulations via the up-to-composition technique. 
%   exhibit two proofs of this result, 
% one in which the  result is directly imported from \piI, the other, presented
% in Appendix~\ref{s:enf},
% that lifts eager normal-form bisimulations into  OGS bisimulations using up-to-composition technique. 

% Moreover,   
% importing   known results in \piI \cite{Adrien} onto OGS, we derive that such
% induced semantics is the equality of 
% Lassen's trees. We also exhibit  a
% direct proof of the result, that does not appeal to  \piI, 
% and that relies  onn the up-to composition technique for OGS.

\item We derive congruence and compositionality properties for the OGS semantics,
as well as a notion of tensor product over configurations that computes
interleavings of traces.
% , exploting 
%  for the mapping into \piI and its algebraic  properties.
\dsOLD{maybe hint at the link between tensor and the compositionality things in the previous sentence?}

\item 
We show that when equality in OGS is based on complete traces,
the semantics induced on $\LasV$ is the same for the 
Alternating, Concurrent, and Well-Bracketed OGS.
\cutLMCSsecondRound{Therefore, drawing from the results about the Well-Bracketed OGS,
such a semantics also  coincides with that produced by contextual equivalence in $\LasV$ extended
with higher-order references.  }

% Again, the results exploits the up-to
% composition technique imported from \piI. We also exploit the mapping into \piI to derive a few
% intermediate OGS properties that have to do with  compositionality.
\cutLMCSsecondRound{
\item We transport the idea of complete traces onto \piI; 
 the formalisation makes use of a standard linear-receptive type system for \piI.
 We thus obtain a new behavioural equivalence for \piI, 
 called \emph{\completetrace equivalence}.  
 The OGS full abstraction result under complete traces 
 can then be transported onto \piI,
 becoming full abstraction for Milner's encoding 
%(under \completetrace equivalence) 
 with respect to contextual equivalence of $\LasV$ extended with higher-order references.

% It is unclear how such a result could be derived in \piI  without going through OGS. 
\dsOLD{probably remove sentence above}}
\end{enumerate} 







% \begin{itemize}
% \item[(i)]
%  transfer  techniques between
% the two models;
% \item[(ii)] 
% switch  freely between the two models in the analysis of 
% the
%  OGS and \piI
%   representations of $\LasV$, 
% so to  establish new results or obtain new proofs.
% \end{itemize} 
% Concering (i), 
% we   transfer the `bisimulation up-to techniques' for  $\piI$, notably a form of `up-to context',
%  onto OGS, yielding a technique called `up-to composition' in OGS.
% The result is a powerful technique that allows us, when needed, to split an OGS
% configuration into more elementary configurations. The technique was inspired by the
% mapping of OGS configurations into \piI and its soundness is derived via mapping itself.
 


% \begin{enumerate}
% \item We exploit such
% a   technique in the proofs of the correspondence between ordinary and concurrent OGS. 
% \item  We transfer the equality produced by OGS complete traces onto the $\pi$-calculus,
%   obtaining a new behavioural equivalence in for \piI.
%  Thus the OGS full abstraction for complete traces in \piI
%  becomes a full abstraction for Milner's encoding wrt contextual
%   equivalence of call-by-value $\lambda$-calculus extended  with higher-order references.
% \item Moreover, in proofs we often  freely switch between the two models..... 
% \finish{give examples} 
% \end{enumerate} 



% A by-product of the study above are a few new developments for OGS.  These include:
% \begin{enumerate}
% \item
% The study of  the concurrent variant of OGS and its correspondence with 
%   ordinary OGS mentioned above;
% \item 
% for the $\LasV$ representation, 
% the proofs that 
%   the  semantics obtained from the \emph{complete traces} (where...), in both
%   ordinary and concurrent OGS,   is the same as that produced by such traces on the 
% WB variant of OGS (where....) . 
% \item 
% a full abstraction result 
% for the complete trace semantics in the  WB OGS 
%  for 
%  contextual
%   equivalence of call-by-value $\lambda$-calculus  with higher-order references.
% \end{enumerate} 




% The tie between the  OGS and $\pi$-calculus  models allow us to transfer results between
% the two models. In the paper this is exploited in several ways, that include the
% following. 

% \begin{enumerate}
% \item We transfer the up-to techniques for the $\pi$-calculus onto OGS, such as forms of
%   `up-to context', yielding a technique called `up-to composition' in OGS. We exploit such
% a   technique in the proofs of the correspondence between ordinary and concurrent OGS. 
% \item  We transfer the equality produced by OGS complete traces onto the $\pi$-calculus,
%   obtaining a new behavioural equivalence in for \piI.
%  Thus the OGS full abstraction for complete traces in \piI
%  becomes a full abstraction for Milner's encoding wrt contextual
%   equivalence of call-by-value $\lambda$-calculus extended  with higher-order references.
% \item Moreover, in proofs we often  freely switch between the two models..... 
% \finish{give examples} 
% \end{enumerate} 

\noindent
The results about OGS in (1-4) are obtained by exploiting
  the mapping  into \piI and  its algebraic  properties and proof techniques, as well as  
the up-to-composition technique for OGS imported from \piI.
Analogous results are established for the call-by-name strategy.
To summarize our results, we consider in this paper the following LTSs and equivalences:
\begin{center}
\begin{tabular}{@{}lccccc@{}}
  \toprule
  LTS & \piI & {\small output-prioritised} \piI & \COGS & \AOGS & \WBOGS \\
  \midrule
  bisimilarity & $\wbPI$ & $\wbOPI$ & $\wbC$ & $\wbA$ & $\wbWB$ \\
  trace equivalence & $\TEPI$ & $\TEN$ & $\TEp$ & $\TEA$ & $\TEwb$ \\
  complete-trace equivalence & & & $\TEpc$ & $\TEc$ & $\TEwbc$ \\
  \bottomrule
\end{tabular}
\end{center}
Terms of the call-by-value $\lambda$-calculus $M,N$ can then be embedded into the LTS $\piI$ and $\opLTS$ with the encoding
$\encoW M$ and into the $\COGS$, $\AOGS$ and $\WBOGS$ as initial configurations $\conf M$.
Using these embeddings, the various equivalence of the previous figure can be lifted into equivalence over terms.
We then prove that:
\begin{enumerate}
\item $\wbPI, \wbOPI, \wbC, \wbA, \wbWB, \TEPI, \TEN, \TEp, \TEA$ all coincides;
\item $\TEpc, \TEc$ and $\TEwbc$ all coincides.
\end{enumerate}
We also prove that the first group of equivalences also coincides with eager normal-form bisimilarity.


\emph{Structure of the paper.} 
Sections~\ref{s:nota} to~\ref{s:encoLpi} contain background material:
general  notations,
 \piI, $\LasV$,
the representations of $\LasV$  in the   Alternating OGS (\AOGS)
and  in \piI.
The following sections contain the new material.
In Section~\ref{s:encoGames} 
we study  the relationship  between the two $\LasV$ representations, 
 in   \piI using the output-prioritised LTS. 
In  Section~\ref{s:cogs}
we establish a similar relationship between 
 a  new  Concurrent OGS (\COGS) and  \piI using its ordinary LTS. We also transport 
up-to techniques onto OGS, and prove that all the semantics of $\LasV$ examined 
(OGS, \piI, traces, bisimulations)
 coincide. 
We import compositionality results for OGS from \piI in Section~\ref{s:compo}.
In Section~\ref{s:wb-ogs} we introduce the Well-Bracketed OGS (\WBOGS) and complete
traces. 
In Section~\ref{sec:tensor},
we introduce a notion of product of OGS configurations, that is used to
provide a wide range of decomposition results.
In  Section~\ref{s:wbc_vs_c} 
we relate all complete-trace equivalences in OGS (\AOGS, \COGS, \WBOGS). 
\cutLMCSsecondRound{In Section~\ref{s:fapiI}  we import   complete-trace equivalence and the full
abstraction result (for pure $\lambda$-terms)  into \piI.}
In Section~\ref{s:cbnNEW} we discuss the case of call-by-name.
Finally, in Section~\ref{s:cf} some further related work and future research directions. 
Appendixes~\ref{a:behav}-\ref{a:cbn} report additional  material for proofs and results in
Sections~\ref{s:encoGames}, \ref{s:cogs}, \ref{s:wb-ogs}, and~\ref{s:cbnNEW}. 

An extended abstract of this paper had appeared at CSL'22
\cite{JaberS22}, without proofs and without the contents in
Sections~\ref{s:wb-ogs}-\ref{s:cbnNEW}. 

% deriving  a full-abstraction result for  an extension of $\La
% contextual equivalence with contexts having store.


% So after presenting notations about LTSs and partial maps in Section~\ref{s:nota},
% we recall \piI, including up-to techniques, in Section~\ref{ss:piI}
% and the call-by-value $\lambda$-calculus $\LasV$ in Section~\ref{ss:cbv}.
% Operational game semantics (OGS) for $\LasV$ (\AOGS) is then introduced in Section~\ref{s:AOGS}.
% The representations of $\LasV$ and OGS into \piI are given in Section~\ref{s:encoLpi}.
% The correspondence between the two models is established 
% using the output-prioritised \opLTS for \piI in Section~\ref{s:encoGames},
% and using a new Concurrent OGS (\COGS), in Section~\ref{s:cogs}.
% There, notion of up-to techniques for OGS are defined, that are used to establish
% the correspondence between \AOGS and \COGS.
% Compositionality results for OGS are established in Section~\ref{s:compo}.
% The Well-Bracketed OGS (\WBOGS) is introduced in Section~\ref{s:wb-ogs},
% together with the notion of complete trace.
% Relating complete trace equivalence for \AOGS, \COGS and \WBOGS is established in 
% Section~\ref{s:wbc_vs_c}.
% Finally, in Section~\ref{s:fapiI} the notion of complete trace equivalence is imported into a new 
% behavioural equivalence for \piI, giving a full-abstraction result wrt
% contextual equivalence with contexts having store.

% \finish{I have cut this: 
% We have mentioned that  game semantics allows one  to obtain a 
%  variety of semantics, corresponding to different discriminating powers that 
% the contexts of the language modelled offer. 
% Something similar can be achieved in the $\pi$-calculus, by means of type systems.
% Intuitively, with a type system, an observer is supposed to follow the discipline of names expressed by
% the types; this may sometimes significantly  reduce the  set of legal observations that
% may be performed on the processes, yiealding coarser forms of process equivalence. 
% Alternatively, explicit constraints  or modifications may be imposed  on the behavioural
% equivalences (e.g., modifying the bisimulation clauses); effectively this means modifying
% the LTS with which an external  observer may interact with the tested processes
% (e.g. \cites{Boreale, SanBook}).  
% }


